% Transcription — fonds Alexandre Grothendieck, shelfmark 76, batch 4 (pages 61-62).
% Established from the facsimile archives/batches/76/batch-04.pdf, cut from a
% copy of 76.pdf fetched through the project's relay
% (https://grothendieck-relay.onrender.com/source/76.pdf); PDF page k =
% archivists' page 60+k, no cover sheet. Per the skill transcribe-grothendieck.
% The folder is sixty-two pages in four batches (1-20, 21-40, 41-60, 61-62),
% transcribed in parallel by four passes; this is the last, two pages long.
% Batches 2 and 3 did not exist when this file was written, so the notation
% of the run it continues (pp. 58-60, cubes C(B, sigma)) was read off the
% facsimile of p. 59-60 only, not off a transcription.
% Pass: Opus 5.5 (claude-opus-5-5), 2026-09-26 — first pass, unchecked against the pages by a human.
%
% Dating and title. The dating is the folder's own, copied verbatim from the
% catalogue; the group « Géométrie et topologie combinatoire » (69-89) holds
% it, and since the folder has a date of its own the group's range is not
% added. The title is the catalogue's, without its final full stop.
%
% Pages skipped: none. Pages summarised: none. Both pages are his hand, black
% felt-tip, fast drafting but mostly legible.
%
% His pagination: none on the leaves. No numbered run in this batch; two
% « Corollaire » open p. 61, following a « Prop » at the foot of p. 60
% (batch 3). No « cf. p. N » cross-reference.
%
% What the batch is about. The end of a run on combinatorial cubes
% C = C(B, sigma) and their antipodal involution sigma. P. 61: two
% corollaries for a cube of dimension 3 (Aut^+(C) ≅ Aut of the 4-element set
% Phi_0(C)/sigma; C ↦ Phi_0(C)/sigma an equivalence from oriented 3-cubes to
% 4-element sets with isomorphisms), then « Cube (combinatoire) gauche »:
% the quotient of the boundary by the antipody is the real projective space
% P(E) of dimension n-1, the facet decomposition passes to the quotient
% (C~ = ∂C^/sigma = ∐ p(F-bar)), each closed facet injects into P(E), and a
% case formula for p(F-bar) ∩ p(F'-bar). P. 62: p(F-bar) ⊂ p(F'-bar) iff
% F-bar ⊂ F'-bar; a small figure (square a,b,c,d → lens with edges a,c and
% b,d); P_R(E) gets a decomposition into topological cells whose pairwise
% intersections are a cell or two disjoint cells, with combinatorial
% structure Phi^{**}/sigma, the opposite of Phi^*/sigma (non-empty partial
% sections of B over I); then « Automorphismes des cubes gauches (n ≥ 3) »:
% the chain Aut(C)/sigma → Aut_proj C~ → Aut_ord(Phi(C~)) and whether the
% composite is surjective; for n odd one recovers (H, sigma) as the oriented
% canonical cover of C~; « Cas n pair : à examiner… » closes the folder.
%
% Notation fixed for the batch, following batch 1: \Phi for his barred-I
% symbol (Phi_0(C), Phi^*, Phi^{**}); \mathbb{B} for his double-stroked B;
% \hat{C} for the C with a small hat or arrow above it; \widetilde{C} for
% his C with a tilde; \mathcal{H} for his script H.
%
% Apparatus: 7 \ill{}, 9 \uncertain{} (plus « sinon », doubtful, flagged in a note). Hardest: the three-line case formula
% on p. 61 (overwritten, third line cancelled) and the two interleaved last
% lines of p. 61.

\documentclass[11pt,a4paper]{article}
% Le préambule commun à toutes les transcriptions du fonds Grothendieck.
%
% Il tient en une idée : **l'incertitude se compose, elle ne se résout pas.**
% Une transcription qui lisse un mot douteux a détruit la seule chose pour
% laquelle on la faisait — un lecteur doit pouvoir savoir, à chaque endroit,
% ce qui était sur le feuillet et ce qui a été ajouté. Les six macros qui
% suivent sont donc l'appareil critique, et non de la décoration.
%
% Les mêmes commandes sont reconnues par `scripts/render.mjs`, qui produit la
% vue de lecture du site. Une macro ajoutée ici doit l'être aussi là-bas,
% faute de quoi le rendu échouera bruyamment — ce qui est voulu : un rendu
% silencieusement faux est pire qu'un rendu absent.

\ProvidesPackage{grothendieck}[2026/08/08 Transcriptions du fonds Grothendieck]

\usepackage{amsmath,amssymb,amsthm}
\usepackage{mathrsfs}
\usepackage{stmaryrd}
% Les diagrammes commutatifs sont partout dans ce fonds : c'est la langue
% même de la géométrie algébrique de Grothendieck, et une transcription qui
% les rendrait en prose ne transcrirait rien.
\usepackage{tikz-cd}
% Français par défaut — c'est la langue des feuillets — mais l'anglais est
% chargé aussi : la traduction et le résumé sont des documents anglais, et la
% typographie française (espaces avant les deux-points) y serait une faute.
% Ces documents basculent par \selectlanguage{english} en tête de corps.
\usepackage[english,french]{babel}
\usepackage{fontspec}
\usepackage{geometry}
\geometry{a4paper,margin=25mm,marginparwidth=28mm}
\usepackage{marginnote}
\usepackage{xcolor}
% ulem, not soul. `soul`'s \st and \ul are built on a character-by-character
% scan that XeLaTeX's font machinery breaks: the document fails with an
% undefined control sequence rather than degrading. ulem does the same two jobs
% and is Unicode-engine safe. `normalem` stops it hijacking \emph, which the
% transcriptions use for Grothendieck's own emphasis.
\usepackage[normalem]{ulem}
\usepackage{hyperref}
\hypersetup{colorlinks=true,linkcolor=black,urlcolor=black,pdfborder={0 0 0}}

% --- Métadonnées du lot -----------------------------------------------------
% Reprises telles quelles par le script de rendu, qui les affiche en tête de
% la vue de lecture. Elles fixent ce que le fichier prétend couvrir : sans
% elles, une transcription flotte sans point d'attache dans un fonds de seize
% mille feuillets.

\newcommand{\folder}[1]{\gdef\grothFolder{#1}}
\newcommand{\batch}[1]{\gdef\grothBatch{#1}}
\newcommand{\leaves}[2]{\gdef\grothFirst{#1}\gdef\grothLast{#2}}
\newcommand{\foldertitle}[1]{\gdef\grothFoldertitle{#1}}
\gdef\grothFolder{?}\gdef\grothBatch{?}\gdef\grothFirst{?}\gdef\grothLast{?}\gdef\grothFoldertitle{}

% --- Le résumé --------------------------------------------------------------
%
% Il ouvre la lecture modernisée, sous le titre « Résumé », et il est composé
% en petit. Ce n'est pas un ornement : c'est le seul passage du document qui
% ne soit pas des mathématiques, et un lecteur doit pouvoir le reconnaître
% avant de l'avoir lu — puis le sauter s'il connaît déjà le sujet, ou s'y
% arrêter s'il ne le connaît pas. Le filet qui le suit marque la reprise du
% corps.

\newenvironment{resume}
  {\small\setlength{\parskip}{0.45em}}
  {\par\medskip\hrule\medskip}

% --- Mots-clés ---------------------------------------------------------------
%
% En anglais, en clôture du résumé de la lecture modernisée : le vocabulaire
% moderne sous lequel chercher ce que les feuillets construisent. Le manifeste
% du site (`npm run manifest`) les extrait pour étiqueter la cote entière —
% cette ligne est la source unique des tags, il n'y a pas de fichier à côté.
\newcommand{\keywords}[1]{\par\smallskip\noindent
  {\footnotesize\textsc{Keywords} --- \textit{#1}}\par}

% --- L'appareil critique ----------------------------------------------------

% Le numéro de feuillet, en marge. C'est l'ancre qui permet de régler un
% désaccord sur une formule en regardant *un* feuillet plutôt qu'une chemise
% de six cents. Le site s'en sert aussi pour tourner les pages du fac-similé
% quand on fait défiler la transcription.
\newcommand{\leaf}[1]{%
  \marginnote{\footnotesize\color{gray!70}\textsf{#1}}%
  \label{leaf:#1}\ignorespaces}

% Illisible : on ne devine pas. Un mot inventé qui se lit comme les autres est
% le pire résultat possible.
\newcommand{\ill}{\textcolor{red!60!black}{[\,\ldots\,]}}

% Lecture douteuse : le mot est donné, et signalé comme incertain.
\newcommand{\uncertain}[1]{\uline{#1}}

% Ajout de l'éditeur — une lettre manquante, un mot sous-entendu.
\newcommand{\add}[1]{\textcolor{blue!50!black}{[#1]}}

% Biffé par Grothendieck lui-même. Ce qu'il a rayé fait partie du document :
% une rature dit souvent où la pensée a changé de direction.
\newcommand{\struck}[1]{\sout{#1}}

% Note du transcripteur, sur l'état matériel du feuillet.
\newcommand{\note}[1]{\par\smallskip{\small\color{gray!60!black}\textit{[#1]}}\par\smallskip}

% Note marginale de Grothendieck, distincte du corps du texte.
\newcommand{\marginal}[1]{\par\smallskip\noindent\hspace{1em}%
  {\small\color{gray!60!black}#1}\par\smallskip}

% --- Titre ------------------------------------------------------------------

\AtBeginDocument{%
  \begin{center}
    {\large\bfseries Fonds Alexandre Grothendieck --- cote n\textsuperscript{o}~\grothFolder}\\[2pt]
    {\itshape\grothFoldertitle}\\[4pt]
    {\small feuillets \grothFirst--\grothLast\ (lot \grothBatch)}\\[2pt]
    {\footnotesize\color{gray!60!black}Transcription établie d'après le fac-similé numérisé
     par l'Université de Montpellier.\\
     Les lectures incertaines sont soulignées, les illisibles notées [\,\ldots\,],
     les ajouts entre crochets.}
  \end{center}
  \vspace{1em}}

\endinput


\folder{76}
\batch{4}
\pages{61}{62}
\dating{[vers 1977]}
\watermark{Édition de démonstration}
\foldertitle{n-cartes cellulaires : notes manuscrites (s.d.)}

\begin{document}

\page{61}
\note{la page fait suite à une « Prop » du bas de la page 60 (les
automorphismes d'un cube induisent $\pm$ l'identité sur les sommets, d'où
$\mathrm{Aut}^+(C) \hookrightarrow \mathrm{Aut}(\Phi_0(C))$ en dimension
impaire)}

\emph{Corollaire} Soit $C$ un cube de dim 3. Alors
\[
  \mathrm{Aut}^+(C) \xrightarrow{\ \sim\ }
  \mathrm{Aut}_{\mathrm{ens}}\bigl(\underbrace{\Phi_0(C)/\sigma}_{\text{ens.\ de card.\ } 4}\bigr)
\]

\emph{Corollaire} Le foncteur $C \mapsto \Phi_0(C)/\sigma$ est une
équivalence de la catégorie des 3-cubes \add{orientés} (\ill{}),
\struck{\uncertain{dans}} $\to$ la cat.\ des ens.\ de card.\ 4 (en prenant
comme morphismes les iso\ldots)
\note{« orientés » est écrit au-dessus de la ligne, souligné, et relié à la
parenthèse qui suit, dont le mot court reste illisible ; le mot biffé en
début de ligne se termine par une flèche}

\emph{Cube (combinatoire) gauche}

\[
  C = C(\mathbb{B}, \sigma) \subset E = E_C = E(\mathbb{B}, \sigma)
\]
cube. Considérons le quotient par la relation d'antipodisme, mais sur
$\partial\hat{C} = \coprod_{F \in \Phi^{**}(C)} \overline{F}$
\note{le $C$ de $\partial\hat{C}$ porte au-dessus un petit signe (chapeau ou
petite flèche), ici et plus bas ; l'argument de $\Phi^{**}$ est surchargé,
lu \uncertain{$C$}}
\[
  \partial\hat{C}/\sigma \simeq E^*/\mathbf{R}^* = P(E)
\]
--- \uncertain{espace} projectif réel de dim $n-1$

La décomposition de $\partial\hat{C}$ en facettes donne une décomposition
(via $C \xrightarrow{\ p\ } \partial C$ \uncertain{passage au quotient})
\note{la flèche est lue telle qu'elle est écrite ; c'est la projection $p$
sur le quotient par $\sigma$ qui sert ensuite}
\[
  \widetilde{C} = \partial\hat{C}/\sigma
  = \coprod_{F \in \Phi^{**}} p(F)
  = \coprod_{F \in \Phi^{**}} p(\overline{F})
\]
\note{le signe entre $\widetilde{C}$ et $\partial\hat{C}$ est surchargé (un
$=$ repassé) ; dans la première somme la lettre devant $(F)$ est corrigée,
lue \uncertain{$p$}}

NB Pour tout $F \in \Phi^{**}$, $\overline{F} \to P(E)$ est injectif.

On a\quad \struck{$p(\overline{F} \cap \overline{F'}) = p(\overline{F}) \cap
p(\overline{F'})$ si $\overline{F} \cap \overline{F'} =$}
\[
  p(\overline{F}) \cap p(\overline{F'}) =
  \left\lbrace
  \begin{array}{ll}
    \varnothing & \text{si } \overline{F} \cap \overline{F'} = \varnothing
      \text{ et } \overline{F} \cap \sigma\overline{F'} = \varnothing \\
    \bigl[p(\overline{F}) \cap p(\overline{F'})\bigr] \cup
      \bigl[p(\overline{F}) \cap p(\overline{F'})\bigr] & \text{sinon}
  \end{array}
  \right.
\]
\note{« sinon » est lu avec doute ; la deuxième ligne est très surchargée : le signe de réunion est
repassé, et dans le second crochet le $\cap$ et le $p$ qui suit sont
récrits, avec un trait au-dessus qui pourrait être un $\sigma$ ou une barre ;
la lecture des deux crochets comme identiques est donc incertaine, et l'on
attendrait $p(\overline{F} \cap \overline{F'}) \cup p(\overline{F} \cap
\sigma\overline{F'})$ ; une troisième ligne, finissant par « $\ne
\varnothing$ », est entièrement biffée et illisible}

(\emph{NB} les relations $\overline{F} \cap \overline{F'} \neq \varnothing$ et
$\overline{F} \cap \sigma\overline{F'} \neq \varnothing$ ne s'excluent pas]
\struck{i.e.\ si $\mathbb{B} \to I$ \uncertain{et} $\mathbb{B}'$,
$\mathbb{B}''$ deux sections partielles}
et $p(\overline{F}) \cap p(\overline{F'})$ \uncertain{n'est pas nécess.}\
connexe ; il est connexe \struck{les} \struck{\ill{}}, si on n'a pas ces 2
relations à la fois)
\note{les deux dernières lignes de la page sont entrelacées : « connexe
\ldots\ à la fois » est écrit en partie sous la ligne qui précède}

\page{62}
\[
  \bigl[p(\overline{F}) \subset p(\overline{F'})\bigr] \Longleftrightarrow
  \overline{F} \subset \overline{F'}
\]

\note{dans la marge gauche, un petit carré dont les côtés portent $a$ (bas),
$b$ (gauche), $c$ (haut), $d$ (droite), et au-dessous une lentille à deux
arcs, l'arc supérieur marqué « $a, c$ », l'inférieur « $b, d$ » : le carré
dont les côtés opposés sont identifiés par l'antipodie}

\emph{NB} On trouve une décomposition de $P_{\mathbf{R}}(E)$ \struck{\ill{}}
en parties qui sont des cellules top.\ --- les intersections entre deux
telles cellules, si elles sont non vides, étant une cellule ou réunion de
2 cellules disjointes (\uncertain{donc ce n'est pas} \ill{} une cellule).
La structure combinatoire de cette division cellulaire (ens.\ ordonné des
parties \ill{}) est \struck{\ill{}} $\Phi^{**}/\sigma \simeq$ ensemble
ordonné opposé de $\Phi^{*}/\sigma$, où $\Phi^{*}$ est l'ens.\ des
sections partielles non vides de $\mathbb{B}$ sur $I$.
\note{l'indice $\mathbf{R}$ de $P_{\mathbf{R}}(E)$ est ajouté après coup}

\emph{Automorphismes des cubes gauches} ($n \geqslant 3$)
\[
  \mathrm{Aut}(C)/\sigma \xrightarrow{\ \sim\ }
  \mathrm{Aut}_{\mathrm{proj}}\,\widetilde{C} \xrightarrow{\ \sim\ }
  \mathrm{Aut}_{\mathrm{ord}}\bigl(\Phi(\widetilde{C})\bigr) \quad ?
\]
\note{la formule est encadrée à gauche par un trait en crochet qui englobe
le titre et la phrase suivante}

On sait déjà que la première flèche est injective, il est évident que la
seconde l'est, le composé \struck{est} \struck{\ill{}} surjectif ?
Si $n$ impair, on récupère $(\mathcal{H}, \sigma)$ \uncertain{avec sa}
division cellulaire, comme le rev.\ orienté canonique de $\widetilde{C}$,
et tt autom.\ de $\Phi(\widetilde{C})$ est induit par un autom.\ orienté de
$\widetilde{C}$.

Cas $n$ pair : à examiner\ldots
\note{fin du dossier}

\end{document}
